\section{Related work}

There is extensive prior work on dynamic load balancing for
shared-memory systems.  The Cilk~\cite{FrigoLeRa98} system was the
first to provide efficient load balancing for a wide variety of
irregular applications.  Load balancing in Cilk applications is
achieved by a scheduler that follows the \textit{depth-first work,
  breadth-first steal} principle~\cite{blumofe:94}. Cilk's scheduling
strategy, in which each worker maintains its own set of
tasks and steals from other workers when it has nothing to do, is
often referred to as \textit{work-stealing}.  
%
Following Cilk, many runtimes for shared memory task parallelism
utilize work-stealing schedulers. 
%
Of these, OpenMP 3.0~\cite{openmp30:08}, Intel's Threading
Building Blocks (TBB)~\cite{reinders:07}, 
Java Fork-Join~\cite{DL2000}, Microsoft's Parallel Patterns
Library (PPL), and Task Parallel Library (TPL) are the most popular.
% Several implementations of X10's async-finish programming model have
utilized work-stealing schedulers including X10's
\textit{breadth-first}~\cite{cong:08} and
\textit{compiler-supported}~\cite{TOWH2012} and Guo et
al.'s~\cite{sarkar:09} hybrid model for work stealing.

The techniques for dynamic load balancing in shared-memory
environments are not directly applicable to distributed-memory
machines because of a variety of issues such as network latency and
bandwidth, and termination detection. 
%
Various bodies of work have addressed the problem of dynamic load balancing on
distributed-memory machines.
% 
Grama et al.~\cite{vipin:94,grama:99} discuss various load balancing 
strategies for distributed parallel searches that are independent of the 
specific search technique.
%
Blumofe et al.~\cite{blumofe:97} adapt the Cilk work-stealing model to 
distributed shared memory by limiting the scope of the programs to be purely 
functional.
%
ATLAS~\cite{baldeschwieler:96} and Satin~\cite{nieuwport:01} both use 
hierarchical work-steal to acheive global load balancing.
%
\Charmpp{}~\cite{sinha:93,kale:93} monitors the execution of its distributed
programs for load imbalance and migrates computation objects to low-load places
to correct the load imbalance.
%
Global load balancing for message passing environments has been researched 
for specific problems by Batoukov and Sorevik~\cite{batoukov:99}.

UTS, an excellent example of an irregular application, was first described by
Prins et al.~\cite{prins:03}.  It has since been widely used as a
benchmark for dynamic load balancing. 
%
Olivier and Prins~\cite{olivier:08} provided the first scalable implementation
of UTS on clusters that provided up to 80\% efficiency on 1024 nodes.
%
To this end, they employed a custom application-level load balancer along with
optimizations such as one-sided communications and novel termination detection
techniques.
% 
Dinan et al.~\cite{dinan:09} provide a framework for
global load balancing, which was used to achieve speedups on UTS on 8196
processors. Global load balancing and termination detection facilities
were provided to express irregular applications.
%
By reserving one core per compute node on the cluster exclusively for lock and
unlock operations, this framework allowed threads to steal work asynchronously
without disrupting the victim threads. However, the cost paid was a static
allocation (one core out of every eight) for communication. This results in
lower throughput because the thread is not available for user-level
computations.
% 
Saraswat et al.~\cite{glb} introduced lifeline-based global load
balancing and showed 87\% efficiency on 2048 nodes. Later work by the
X10 team demonstrated 98\% parallel efficiency with 55,680 Power7
cores~\cite{tardieu-petascale}.  An implementation of the life-line
algorithm in Co-Array Fortran achieved 58\% efficiency at 8192
nodes~\cite{CAF-HPCC2011}. A more recent UTS code using CAF 2.0 {\tt
  finish} construct achieves a 74\% parallel efficiency on 32,768
Jaguar cores~\cite{Chaoran:2013}.

%% Dave: Leaving this point out since we're tight on space and short
%%  on time...
%
% More broadly, the approach to programming taken by our GLB library
% in which the end-user of a framework provides snippets of
%sequential code that mesh with an underlying distributed framework is
% a commonly applied paradigm. Systems such as J2EE and MapReduce...


